\honest{The Genetic Fallacy}{Eliezer Yudkowsky}{2008}

Lists of fallacies include the genetic fallacy: attacking a belief because of why someone holds it. At first sight this is ``a very strange idea.'' If the causes of a belief do not make it reliable, what does? We trust Deep Blue's chess moves because we understand its code.

The articles say origins sometimes matter, as with a trusted expert, and sometimes do not, as with Kekulé, who first saw the benzene ring in a dream. ``So sometimes the genetic fallacy is a fallacy, and sometimes it's not?'' I ask. \nb{By the end of the post Yudkowsky reaches the same position: the source stops mattering once there is clear-cut evidence, and deserves attention when there is none.}

It is a fallacy in principle, because a belief's original cause is not its current justification. But we change our minds less often than we think, so origins have more force among humans than among ideal Bayesians. When a source you relied on comes under suspicion, clear your mind. \nb{The job-choice survey behind ``We Change Our Minds Less Often Than We Think'' does not concern discredited sources. The belief-perseverance studies do, and support the point.}

Take people who stopped believing that God wrote the Bible but still think it full of ``indispensable ethical wisdom.'' ``They have failed to clear their minds,'' I say. \nb{No case is given. Whether a given person re-examined the Bible's ethics is a question of fact about that person.} They could do significantly better by doubting anything the Bible said because the Bible said it, without simply negating it, since reversed stupidity is not intelligence.

An idea finds support everywhere once it is in your head, so when its source falls, suspect every leaf on that branch. That is hard: ``It takes a convulsive effort to actually reconsider.'' \nb{In the first belief-perseverance study, a debriefing that explained the perseverance process to subjects ``generally proved sufficient to eliminate erroneous self-perceptions.'' That was one laboratory task.} If all the ideas from a source you now distrust still seem right, be ``extremely suspicious.''

With enough clear-cut evidence, as for the benzene ring, the source no longer matters. Without it, give experts more credence and distrust suspect motives. The genetic fallacy is a fallacy when other justifications exist but the genetic accusation ``is presented as if it settled the issue.'' Hal Finney suggests calling correct appeals to origins ``the genetic heuristic.'' \nb{That sentence was added after the post appeared, from a comment Finney left on it.}

I end with four rules of thumb. Distrust genetic accusations against beliefs you dislike. Prefer technical arguments to psychoanalysis. Doubt everything that grew from a discredited source, without rejecting it outright. And be very suspicious if you still believe the early suggestions of a source you rejected.
